%------------------------------------------------------------------
% Questão 4.6 – Velocidade máxima de um animal por análise dimensional
%------------------------------------------------------------------

\textbf{Passo 1 – Massa em função do tamanho}

Todos os animais têm densidade similar (próxima à da água). Logo, a massa é proporcional ao volume:
\[
m = \rho L^3,
\]
onde $L$ é o tamanho característico (comprimento) e $\rho \approx 1000\,\text{kg/m}^3$.

\textbf{Passo 2 – Força muscular máxima}

Um músculo pode exercer força proporcional à sua área de seção transversal:
\[
F_{\max} = \sigma_m A_m,
\]
onde $\sigma_m$ é a tensão máxima do tecido muscular (constante para mamíferos) e $A_m \sim L^2$ (área cresce com $L^2$).

Portanto: $F_{\max} = \sigma_m L^2$.

\textbf{Passo 3 – Trabalho por passada}

Uma passada tem comprimento $\ell \sim L$. O trabalho realizado é:
\[
W = F_{\max} \cdot \ell = \sigma_m L^2 \cdot L = \sigma_m L^3.
\]

\textbf{Passo 4 – Conversão em energia cinética}

Esse trabalho vira velocidade máxima do animal:
\[
W = \frac{1}{2}m v_{\max}^2
\implies
\sigma_m L^3 = \frac{1}{2}\rho L^3 v_{\max}^2.
\]

\textbf{Passo 5 – Cancelar $L^3$ e resolver}

\[
\sigma_m = \frac{1}{2}\rho v_{\max}^2
\implies
v_{\max} = \sqrt{\frac{2\sigma_m}{\rho}}.
\]

\textbf{Resultado importante:}

A velocidade máxima \emph{não depende do tamanho} $L$! Depende apenas de propriedades do material.

\textbf{Verificação dimensional:}
\[
[v_{\max}] = \sqrt{\frac{\text{Pa}}{\text{kg/m}^3}} = \sqrt{\frac{\text{m/s}^2}{\text{m/s}^2 \cdot \text{m}^{-3}}} = \sqrt{\text{m}^2/\text{s}^2} = \text{m/s}\quad\checkmark
\]

\textbf{Aplicação numérica:}

Para mamíferos: $\sigma_m \approx 2\times10^5\,\text{Pa}$ e $\rho \approx 1000\,\text{kg/m}^3$.

\[
v_{\max} = \sqrt{\frac{2 \times 2\times10^5}{1000}} = \sqrt{400} = 20\,\text{m/s} = 72\,\text{km/h}.
\]

Este valor concorda com velocidades máximas observadas na natureza (cavalos, guepardos).

\textbf{Conclusão:} Um rato e um cavalo têm aproximadamente a mesma velocidade máxima porque ambos dependem das mesmas propriedades do tecido muscular!

